% =====================================================================
%  CARNET D'ENTRAÎNEMENT — Fiche 12 (Chimie organique)
%  THÈME 2 — Reconnaissance et classification des groupements fonctionnels
%  NIVEAU DIFFICILE — ÉNONCÉS
% =====================================================================
\documentclass[11pt,a4paper]{article}
\usepackage[utf8]{inputenc}
\usepackage[T1]{fontenc}
\usepackage{lmodern}
\usepackage[french]{babel}
\usepackage{amsmath,amssymb,mathtools}
\providecommand{\degree}{^\circ}
\usepackage{icomma}
\usepackage{siunitx}
\sisetup{output-decimal-marker={,},per-mode=symbol}
\usepackage[version=4]{mhchem}
\usepackage{chemfig}
\setchemfig{atom sep=2em,bond length=0.9em,cram width=2.5pt}
\usepackage{xcolor}
\usepackage{array}
\usepackage{booktabs}
\usepackage{tabularx}
\usepackage[most]{tcolorbox}
\usepackage{enumitem}
\usepackage{tikz}
\usetikzlibrary{arrows.meta,positioning,calc}
\usepackage{fancyhdr}
\usepackage{titlesec}
\usepackage[hidelinks]{hyperref}
\usepackage[a4paper,margin=2.1cm,headheight=26pt,headsep=9pt]{geometry}

\definecolor{primary}{RGB}{150,98,162}
\definecolor{primarydark}{RGB}{92,52,110}
\definecolor{excol}{RGB}{0,135,90}
\definecolor{attncol}{RGB}{200,40,55}
\definecolor{methcol}{RGB}{90,90,100}

\newcommand{\runtitle}{Thème 2 — Difficile — Énoncés}
\pagestyle{fancy}\fancyhf{}
\fancyhead[L]{\small\textcolor{primarydark}{\textbf{Fiche 12} — Chimie organique}}
\fancyhead[R]{\small\textcolor{primarydark}{\runtitle}}
\fancyfoot[C]{\small\thepage}
\renewcommand{\headrulewidth}{0.8pt}
\renewcommand{\headrule}{\hbox to\headwidth{\color{primary}\leaders\hrule height \headrulewidth\hfill}}
\setlength{\parindent}{0pt}\setlength{\parskip}{3pt}

\tcbset{boxbase/.style={enhanced,breakable,boxrule=0.9pt,arc=2.5pt,
   left=3mm,right=3mm,top=2mm,bottom=2mm,fonttitle=\bfseries,coltitle=white,
   toptitle=0.5mm,bottomtitle=0.5mm}}
\newtcolorbox[auto counter]{exo}[1]{boxbase,colback=primary!4,colframe=primary,
   title={\textsc{Exercice}~\thetcbcounter\ \textmd{--- \itshape #1}}}
\newtcolorbox{consigne}{boxbase,colback=methcol!8,colframe=methcol,sharp corners,
   title={\textsc{Consignes}}}

\newcommand{\banner}[2]{%
\begin{tcolorbox}[enhanced,colback=primary,colframe=primary,arc=5pt,boxrule=0pt,
   left=5mm,right=5mm,top=2.5mm,bottom=2.5mm,coltext=white]
{\large\bfseries #1}\\[1pt]{\small #2}
\end{tcolorbox}\vspace{2pt}}

\begin{document}

\banner{Thème 2 — Reconnaissance des groupements fonctionnels}{Niveau \textbf{Difficile} — Énoncés \quad$\vert$\quad Fiche 12, Chimie organique}

\begin{consigne}
Analyse fonctionnelle de \textbf{molécules réelles} (médicaments, biomolécules), comme au concours. Sous-questions progressives. Rappel : un $\ce{-OH}$ sur un carbone \emph{aromatique} (cycle benzénique) s'appelle un \textbf{phénol}, à distinguer de l'alcool (dont le C--OH est \emph{saturé}).
\end{consigne}

% ---------------------------------------------------------------------
\begin{exo}{l'aspirine (acide acétylsalicylique)}
La structure de l'\textbf{aspirine} est représentée ci-dessous :
\begin{center}\chemfig{*6(-=-(-(-[:30]OH)=[:90]O)=-(-O-[:150](-[:-150]CH_3)=[:90]O)=-=)}\end{center}
\begin{enumerate}[label=\textbf{\alph*)},leftmargin=2em,itemsep=1pt]
  \item Recense \textbf{toutes} les fonctions chimiques de la molécule.
  \item On lit parfois « l'aspirine contient une fonction éther » à cause de l'enchaînement $\ce{Ar-O-C}$. Explique pourquoi c'est \textbf{faux}.
  \item Combien de doubles liaisons $\ce{C=C}$ le noyau aromatique apporte-t-il ?
\end{enumerate}
\end{exo}

% ---------------------------------------------------------------------
\begin{exo}{le paracétamol}
Le \textbf{paracétamol} (antalgique) a la structure suivante :
\begin{center}\chemfig{*6(=-=(-NH-[:30](-[:30]CH_3)=[:90]O)-=(-OH)-)}\end{center}
\begin{enumerate}[label=\textbf{\alph*)},leftmargin=2em,itemsep=1pt]
  \item Deux groupements sont greffés sur le cycle. Le motif contenant l'azote est-il une \textbf{amine} ou une \textbf{amide} ? Justifie.
  \item Le groupe $\ce{-OH}$ porté par le cycle est-il une fonction \textbf{alcool} au sens strict de la fiche ? Comment l'appelle-t-on ?
  \item Combien de doubles liaisons $\ce{C=C}$ le noyau apporte-t-il ?
\end{enumerate}
\end{exo}

% ---------------------------------------------------------------------
\begin{exo}{molécule polyfonctionnelle oxygénée (type stéroïde)}
Une biomolécule (fragment simplifié d'un stéroïde du type hydrocortisone) présente le squelette oxygéné suivant :
\begin{center}\chemfig{HO-[:30]CH_2-[:-30](=[:-90]O)-[:30]CH_2-[:-30]CH(-[:-90]OH)-[:30]CH_2-[:-30](=[:-90]O)-[:30]C(-[:90]CH_3)(-[:150]CH_3)-[:30]OH}\end{center}
\begin{enumerate}[label=\textbf{\alph*)},leftmargin=2em,itemsep=1pt]
  \item Combien de fonctions \textbf{cétone} contient la molécule ? Justifie par le voisinage de chaque $\ce{C=O}$.
  \item Combien de fonctions \textbf{alcool} ? Classe \emph{chacune} (I / II / III).
  \item Réponds à la question type concours : « quelles sont les fonctions oxygénées présentes ? » (donne \emph{nombre} et \emph{classe}).
\end{enumerate}
\end{exo}

% ---------------------------------------------------------------------
\begin{exo}{un fragment de type aspartame (acide + amine + amide + ester)}
La molécule ci-dessous reprend le motif d'un édulcorant de type \textbf{aspartame} :
\begin{center}\chemfig{HO-[:30](=[:90]O)-[:-30]CH_2-[:30]CH(-[:90]NH_2)-[:-30](=[:-90]O)-[:30]NH-[:-30]CH_2-[:30](=[:90]O)-[:-30]O-[:30]CH_3}\end{center}
\begin{enumerate}[label=\textbf{\alph*)},leftmargin=2em,itemsep=1pt]
  \item Recense les \textbf{quatre} fonctions présentes (parcours la chaîne, repère d'abord les $\ce{C=O}$).
  \item Le groupe $\ce{-NH2}$ est proche d'un carbonyle. Est-ce une amine ou une amide ? Justifie et donne sa classe.
  \item Deux fonctions portent un $\ce{C=O}$ : comment distingue-t-on l'\textbf{amide} de l'\textbf{ester} ?
\end{enumerate}
\end{exo}

\end{document}
